\chapter{Introducción}
El presente proyecto tiene como objetivo central el diseño y modelado de un sistema integral
para la gestión del Concurso de Personal, abarcando todas las fases críticas del proceso de
reclutamiento: desde la convocatoria inicial y la recepción de postulaciones, pasando por la
evaluación y selección, hasta la contratación efectiva del colaborador. En un contexto laboral
cada vez más competitivo, donde la atracción y retención de talento se convierten en ventajas
estratégicas clave, este estudio se centra en la empresa Tambo, líder indiscutible en el sector
de tiendas de conveniencia en el Perú.Con una estrategia de expansión masiva que ha posicio
nado más de 500 locales a nivel nacional, Tambo enfrenta el desafío de procesos de captación
de talento que deben ser no solo ágiles y precisos, sino también escalables para responder a
un volumen creciente de vacantes. Este proyecto busca optimizar el flujo de información y la
toma de decisiones mediante una herramienta tecnológica innovadora que automatice tareas
repetitivas como el triaje de currículos y el seguimiento de candidatos y centralice la base de
datos de postulantes en una plataforma unificada. De esta manera, se alinea con las demandas
del mercado actual, caracterizado por la agilidad digital y la precisión en la selección de perfiles
idóneos, contribuyendo directamente a la eficiencia operativa y al crecimiento sostenible de la
organización.